%===============================================================================
% LaTeX sjabloon voor de bachelorproef toegepaste informatica aan HOGENT
% Meer info op https://github.com/HoGentTIN/latex-hogent-report
%===============================================================================

\documentclass[dutch,dit,thesis]{hogentreport}

% TODO:
% - If necessary, replace the option `dit`' with your own department!
%   Valid entries are dbo, dbt, dgz, dit, dlo, dog, dsa, soa
% - If you write your thesis in English (remark: only possible after getting
%   explicit approval!), remove the option "dutch," or replace with "english".

\usepackage{lipsum} % For blind text, can be removed after adding actual content

%% Pictures to include in the text can be put in the graphics/ folder
\graphicspath{{../graphics/}}

%% For source code highlighting, requires pygments to be installed
%% Compile with the -shell-escape flag!
%% \usepackage[chapter]{minted}
%% If you compile with the make_thesis.{bat,sh} script, use the following
%% import instead:
\usepackage[chapter,outputdir=../output]{minted}
\usemintedstyle{solarized-light}

%% Formatting for minted environments.
\setminted{%
    autogobble,
    frame=lines,
    breaklines,
    linenos,
    tabsize=4
}

%% Ensure the list of listings is in the table of contents
\renewcommand\listoflistingscaption{%
    \IfLanguageName{dutch}{Lijst van codefragmenten}{List of listings}
}
\renewcommand\listingscaption{%
    \IfLanguageName{dutch}{Codefragment}{Listing}
}
\renewcommand*\listoflistings{%
    \cleardoublepage\phantomsection\addcontentsline{toc}{chapter}{\listoflistingscaption}%
    \listof{listing}{\listoflistingscaption}%
}

% Other packages not already included can be imported here

%%---------- Document metadata -------------------------------------------------
% TODO: Replace this with your own information
\author{Arthur Van Oudenhove}
\supervisor{Dhr. L. Smits}
\cosupervisor{Dhr. A. Algaba}
\title[Optionele ondertitel]%
    {Beveiliging van AI‑agents in Fintech: automatische detectie van beveiligingslekken in integraties op basis van het Model Context Protocol (MCP)}
\academicyear{\advance\year by -1 \the\year--\advance\year by 1 \the\year}
\examperiod{1}
\degreesought{\IfLanguageName{dutch}{Professionele bachelor in de toegepaste informatica}{Bachelor of applied computer science}}
\partialthesis{false} %% To display 'in partial fulfilment'
%\institution{Internshipcompany BVBA.}

%% Add global exceptions to the hyphenation here
\hyphenation{back-slash}

%% The bibliography (style and settings are  found in hogentthesis.cls)
\addbibresource{bachproef.bib}            %% Bibliography file
\addbibresource{../voorstel/voorstel.bib} %% Bibliography research proposal
\defbibheading{bibempty}{}

%% Prevent empty pages for right-handed chapter starts in twoside mode
\renewcommand{\cleardoublepage}{\clearpage}

\renewcommand{\arraystretch}{1.2}

%% Content starts here.
\begin{document}

%---------- Front matter -------------------------------------------------------

\frontmatter

\hypersetup{pageanchor=false} %% Disable page numbering references
%% Render a Dutch outer title page if the main language is English
\IfLanguageName{english}{%
    %% If necessary, information can be changed here
    \degreesought{Professionele Bachelor toegepaste informatica}%
    \begin{otherlanguage}{dutch}%
       \maketitle%
    \end{otherlanguage}%
}{}

%% Generates title page content
\maketitle
\hypersetup{pageanchor=true}

%%=============================================================================
%% Voorwoord
%%=============================================================================

\chapter*{\IfLanguageName{dutch}{Woord vooraf}{Preface}}%
\label{ch:voorwoord}

Mijn interesse in artificiële intelligentie heeft me al vroeg doen zoeken naar de protocollen en standaarden die AI-agents in de praktijk bruikbaar maken. Zo stuitte ik op het Model Context Protocol, een open specificatie die nog geen jaar oud was toen ik eraan begon. Juist die jonge leeftijd trok me: een protocol dat breed wordt omarmd voordat de beveiligingsimplicaties goed begrepen zijn, is precies het soort blinde vlek dat je als onderzoeker wil aanpakken. Dat werd het vertrekpunt voor deze bachelorproef.

Graag dank ik mijn promotor Lieven Smits en mijn co-promotoren Andres Algaba en Brecht Verbeken voor hun begeleiding doorheen dit traject. Hun feedback heeft me geholpen om het onderzoek scherp te houden en de juiste keuzes te maken.
%%=============================================================================
%% Samenvatting
%%=============================================================================

% TODO: De "abstract" of samenvatting is een kernachtige (~ 1 blz. voor een
% thesis) synthese van het document.
%
% Een goede abstract biedt een kernachtig antwoord op volgende vragen:
%
% 1. Waarover gaat de bachelorproef?
% 2. Waarom heb je er over geschreven?
% 3. Hoe heb je het onderzoek uitgevoerd?
% 4. Wat waren de resultaten? Wat blijkt uit je onderzoek?
% 5. Wat betekenen je resultaten? Wat is de relevantie voor het werkveld?
%
% Daarom bestaat een abstract uit volgende componenten:
%
% - inleiding + kaderen thema
% - probleemstelling
% - (centrale) onderzoeksvraag
% - onderzoeksdoelstelling
% - methodologie
% - resultaten (beperk tot de belangrijkste, relevant voor de onderzoeksvraag)
% - conclusies, aanbevelingen, beperkingen
%
% LET OP! Een samenvatting is GEEN voorwoord!

%%---------- Nederlandse samenvatting -----------------------------------------
%
% TODO: Als je je bachelorproef in het Engels schrijft, moet je eerst een
% Nederlandse samenvatting invoegen. Haal daarvoor onderstaande code uit
% commentaar.
% Wie zijn bachelorproef in het Nederlands schrijft, kan dit negeren, de inhoud
% wordt niet in het document ingevoegd.

\IfLanguageName{english}{%
\selectlanguage{dutch}
\chapter*{Samenvatting}
\lipsum[1-4]
\selectlanguage{english}
}{}

%%---------- Samenvatting -----------------------------------------------------
% De samenvatting in de hoofdtaal van het document

\chapter*{\IfLanguageName{dutch}{Samenvatting}{Abstract}}

Autonome AI-agents verwerken in de Fintech-sector steeds vaker betalingstransacties en andere bedrijfskritische taken. De communicatie met interne systemen verloopt daarbij via het Model Context Protocol (MCP), een standaard die integratie vergemakkelijkt maar tegelijkertijd nieuwe aanvalsvlakken opent: \emph{prompt injection}, \emph{tool poisoning} en \emph{privilege escalation}. Generieke beveiligingstools zoals Burp Suite en Semgrep missen de contextuele kennis om MCP-specifieke kwetsbaarheden systematisch op te sporen, en handmatige code reviews zijn arbeidsintensief en foutgevoelig. Dat 56\% van de bedrijven die AI-agents inzetten beveiligingsincidenten rapporteert in hun MCP-omgevingen, onderstreept de nood aan gespecialiseerde tooling.

De centrale vraag is of een geautomatiseerd security assessment framework MCP-server\-implementaties systematisch kan testen op kritieke kwetsbaarheden, en in hoeverre de effectiviteit verschilt van een gestructureerde handmatige beoordeling. Als evaluatiecase dient een opzettelijk kwetsbare testserver (\emph{Goat}) met vijf bekende, opzettelijk ingebedde kwetsbaarheden.

Het framework combineert statische analyse (SAST) met dynamische fuzzing en draait containergebaseerd via Docker~Compose. Dezelfde Goat-server werd parallel handmatig beoordeeld via code review en HTTP-tests als kwantitatieve benchmark.

Het framework detecteerde alle vijf ingebedde kwetsbaarheden in minder dan 30~seconden, bij een precisie van minimaal 92{,}3\% (12 terechte bevindingen op 13 totaal; één bevinding is debateerbaar als false positive). De handmatige procedure vereiste 47~minuten en identificeerde 4 van de vijf kwetsbaarheden.

De resultaten bevestigen dat geautomatiseerde MCP security assessment haalbaar is en in deze evaluatiecase handmatige beoordeling ruim overtreft in snelheid, bij vergelijkbare of betere recall. De externe validiteit blijft beperkt tot één gecontroleerde case met één reviewer; bredere uitspraken over MCP-servers in het algemeen zijn niet gerechtvaardigd. De voornaamste aanbevelingen zijn: autorisatiechecks verplichten per tool, tool outputs sanitiseren als onvertrouwde invoer, het framework inzetten als screeningslaag en de scanner integreren in CI/CD voor continue compliancemonitoring.


%---------- Inhoud, lijst figuren, ... -----------------------------------------

\tableofcontents

% In a list of figures, the complete caption will be included. To prevent this,
% ALWAYS add a short description in the caption!
%
%  \caption[short description]{elaborate description}
%
% If you do, only the short description will be used in the list of figures

\listoffigures

% If you included tables and/or source code listings, uncomment the appropriate
% lines.
\listoftables

\listoflistings

% Als je een lijst van afkortingen of termen wil toevoegen, dan hoort die
% hier thuis. Gebruik bijvoorbeeld de ``glossaries'' package.
% https://www.overleaf.com/learn/latex/Glossaries

%---------- Kern ---------------------------------------------------------------

\mainmatter{}

% De eerste hoofdstukken van een bachelorproef zijn meestal een inleiding op
% het onderwerp, literatuurstudie en verantwoording methodologie.
% Aarzel niet om een meer beschrijvende titel aan deze hoofdstukken te geven of
% om bijvoorbeeld de inleiding en/of stand van zaken over meerdere hoofdstukken
% te verspreiden!

%%=============================================================================
%% Inleiding
%%=============================================================================

\chapter{\IfLanguageName{dutch}{Inleiding}{Introduction}}%
\label{ch:inleiding}


De sector van financiële technologie, of Fintech, heeft de afgelopen vijf jaar een jaarlijkse groei van 23,58\% doorgemaakt \autocite{statista2024fintechgrowth}. Een belangrijk deel van die groei is toe te schrijven aan de inzet van autonome AI-agents: softwaresystemen die zelfstandig beslissingen nemen en acties uitvoeren \autocite{russell2021artificialintelligence}. Ze worden niet alleen voor eenvoudige taken gebruikt, maar ook voor complexe bedrijfsprocessen zoals het afhandelen van klantvragen of het verwerken van transacties. Om te communiceren met interne databases en API's is het Model Context Protocol (MCP) de voorbije jaren de standaard geworden. MCP biedt een uniforme interface waardoor AI-agents platformonafhankelijk met bedrijfssystemen kunnen communiceren, vergelijkbaar met de rol die HTTP speelde voor webapplicaties \autocite{anthropic2024mcpwhitepaper}. Zonder zo'n gestandaardiseerd protocol zou elke organisatie voor elk AI-model een aparte integratie moeten bouwen, met alle gevolgen voor onderhoud en schaalbaarheid van dien.

Die standaardisatie heeft echter ook een keerzijde op vlak van beveiliging. Neem een Fintech-bedrijf dat via MCP een AI-agent toegang geeft tot zijn betaalsystemen. Aanvallers kunnen in zo'n opstelling \emph{prompt injection} gebruiken om de instructies van de agent te manipuleren door gewijzigde invoer in het systeem te injecteren. Via \emph{tool poisoning} kunnen zij schadelijke commando's verbergen in schijnbaar legitieme datastromen. En via \emph{privilege escalation} kan een agent door een beveiligingsfout onbedoeld verhoogde toegangsrechten verkrijgen, waardoor kritieke systemen bereikbaar worden voor onbevoegden \autocite{song2025protocolunveilingattackvectors}.

Security-teams hebben doorgaans te weinig middelen om MCP-implementaties grondig te beveiligen tegen deze specifieke dreigingen. De huidige praktijk is afhankelijk van handmatige code-reviews en ad-hoc beveiligingstools die niet specifiek zijn ontworpen voor de technische eigenheid van het Model Context Protocol. Deze algemene beveiligingstools (zoals Burp Suite voor API-testing of Semgrep voor code-analyse) missen de nodige contextuele kennis om MCP-specifieke kwetsbaarheden, zoals gemanipuleerde tool-\hspace{0pt}beschrijvingen of onveilige JSON-RPC-communicatie, systematisch te detecteren. Dit gebrek aan gespecialiseerde tooling creëert een bottleneck in het ontwikkelproces: ofwel vertraagt innovatie door langdurige handmatige beveiligingscontroles, ofwel worden MCP-servers in productie gebracht zonder grondige veiligheidsbeoordeling. Dat 56\% van de ondernemingen die AI-agents inzetten beveiligingsincidenten rapporteert in hun MCP-omgevingen maakt de urgentie concreet \autocite{guo2025systematicanalysismcpsecurity}. Dit onderzoek bouwt een security assessment framework dat MCP-servers automatisch en systematisch test.


%De inleiding moet de lezer net genoeg informatie verschaffen om het onderwerp te begrijpen en in te zien waarom de onderzoeksvraag de moeite waard is om te onderzoeken. In de inleiding ga je literatuurverwijzingen beperken, zodat de tekst vlot leesbaar blijft. Je kan de inleiding verder onderverdelen in secties als dit de tekst verduidelijkt. Zaken die aan bod kunnen komen in de inleiding~\autocite{Pollefliet2011}:

%\begin{itemize}
%  \item context, achtergrond
%  \item afbakenen van het onderwerp
%  \item verantwoording van het onderwerp, methodologie
%  \item probleemstelling
%  \item onderzoeksdoelstelling
%  \item onderzoeksvraag
%  \item \ldots
%\end{itemize}

\section{\IfLanguageName{dutch}{Probleemstelling}{Problem Statement}}%
\label{sec:probleemstelling}

Fintech-bedrijven beschikken momenteel niet over een gestructureerde en geautomatiseerde manier om MCP-servers op kritieke veiligheidskwetsbaarheden te controleren vóór live-gang. Security teams vertrouwen op handmatige code reviews, een aanpak die arbeidsintensief en foutgevoelig is. Ad-hoc tools zoals Burp Suite of Semgrep bieden slechts gedeeltelijke dekking omdat ze de specifieke context van MCP-aanvallen missen.

\begin{itemize}
  \item \textbf{Specifieke risico's:} Zonder degelijke screening zijn servers kwetsbaar voor:
        \begin{itemize}
          \item \emph{Prompt injection attacks} via manipulatie van tool-beschrijvingen.
          \item \emph{Token theft \& privilege escalation}, waarbij agents te brede rechten krijgen.
          \item \emph{Tool poisoning}, waarbij kwaadaardige instructies in data worden verstopt.
          \item \emph{Confused deputy attacks}, omdat servers vaak geen onderscheid maken tussen gebruikers.
          \item Gebrek aan \emph{audit logging}, waardoor acties niet traceerbaar zijn.
        \end{itemize}
  \item \textbf{Doelgroep:} Het onderzoek richt zich op security teams, DevSecOps-engineers en enterprise architects binnen Fintech, e-commerce en andere gereguleerde industrieën.
\end{itemize}

%Uit je probleemstelling moet duidelijk zijn dat je onderzoek een meerwaarde heeft voor een concrete doelgroep. De doelgroep moet goed gedefinieerd en afgelijnd zijn. Doelgroepen als ``bedrijven,'' ``KMO's'', systeembeheerders, enz.~zijn nog te vaag. Als je een lijstje kan maken van de personen/organisaties die een meerwaarde zullen vinden in deze bachelorproef (dit is eigenlijk je steekproefkader), dan is dat een indicatie dat de doelgroep goed gedefinieerd is. Dit kan een enkel bedrijf zijn of zelfs één persoon (je co-promotor/opdrachtgever).

\section{\IfLanguageName{dutch}{Onderzoeksvraag}{Research question}}%
\label{sec:onderzoeksvraag}

\paragraph{Hoofdvraag}
Hoe kan een framework worden ontwikkeld dat MCP-server implementaties automatisch test op een vooraf gedefinieerde set kritieke kwetsbaarheden, en hoe verhoudt de dekkingsgraad en benodigde tijd per assessment zich tot een handmatige beoordeling op basis van twee concrete evaluatiecases?

\paragraph{Deelvragen Probleemdomein}
\begin{enumerate}
  \item Welke kwetsbaarheden uit de OWASP MCP Top 10 zijn het meest kritiek voor Fintech-toepassingen, en hoe worden ze geclassificeerd in de bestaande literatuur?
  \item Welke beperkingen hebben generieke beveiligingstools zoals Burp Suite en Semgrep bij het detecteren van MCP-specifieke kwetsbaarheden?
  \item Hoeveel tijd vereist een handmatige beoordeling van een MCP-server implementatie volgens beschikbare richtlijnen, en wat zijn de voornaamste bronnen van foutmarge?
\end{enumerate}

\paragraph{Deelvragen Oplossingsdomein}
\begin{enumerate}
  \item Welk percentage van de ingebedde kwetsbaarheden in de evaluatiecases detecteert het framework correct (recall), en wat is de bijhorende false positive rate?
  \item Hoeveel tijd heeft het framework nodig voor een volledige assessment in vergelijking met een gestructureerde handmatige procedure op dezelfde target?
  \item Welke categorieën kwetsbaarheden detecteert het framework betrouwbaarder dan handmatige review, en in welke gevallen blijft menselijke expertise noodzakelijk?
\end{enumerate}

%Wees zo concreet mogelijk bij het formuleren van je onderzoeksvraag. Een onderzoeksvraag is trouwens iets waar nog niemand op dit moment een antwoord heeft (voor zover je kan nagaan). Het opzoeken van bestaande informatie (bv. ``welke tools bestaan er voor deze toepassing?'') is dus geen onderzoeksvraag. Je kan de onderzoeksvraag verder specifiëren in deelvragen. Bv.~als je onderzoek gaat over performantiemetingen, dan

\section{\IfLanguageName{dutch}{Onderzoeksdoelstelling}{Research objective}}%
\label{sec:onderzoeksdoelstelling}

Deze bachelorproef beoogt een geautomatiseerd security assessment framework te ontwerpen, ontwikkelen en evalueren dat specifiek is toegespitst op MCP-server implementaties. De evaluatie steunt op twee concrete cases: een opzettelijk kwetsbare testserver (\emph{Goat}) en een representatieve open-source MCP-server. De concrete deliverables zijn:
\begin{enumerate}
  \item Een werkend prototype (Proof-of-Concept) van de security scanner die kwetsbaarheden uit de OWASP MCP Top 10 kan detecteren via statische analyse en dynamische fuzzing.
  \item Een kwantitatieve evaluatie van de effectiviteit van het framework op basis van twee cases, gemeten aan de hand van:
        \begin{itemize}
          \item Detectiegraad (recall): het percentage correct geïdentificeerde kwetsbaarheden
          \item Tijdsverschil: benodigde tijd voor een volledige assessment tegenover een gestructureerde handmatige procedure
          \item False positive rate: nauwkeurigheid van de detectie
        \end{itemize}
  \item Een set evidence-based aanbevelingen voor de veilige implementatie van MCP, onderbouwd door de bevindingen uit de evaluatiefase, met expliciete vermelding van de beperkingen van de externe validiteit.
\end{enumerate}


%Wat is het beoogde resultaat van je bachelorproef? Wat zijn de criteria voor succes? Beschrijf die zo concreet mogelijk. Gaat het bv.\ om een proof-of-concept, een prototype, een verslag met aanbevelingen, een vergelijkende studie, enz.

\section{\IfLanguageName{dutch}{Opzet van deze bachelorproef}{Structure of this bachelor thesis}}%
\label{sec:opzet-bachelorproef}

% Het is gebruikelijk aan het einde van de inleiding een overzicht te
% geven van de opbouw van de rest van de tekst. Deze sectie bevat al een aanzet
% die je kan aanvullen/aanpassen in functie van je eigen tekst.

De rest van deze bachelorproef is als volgt opgebouwd:

In Hoofdstuk~\ref{ch:stand-van-zaken} wordt een overzicht gegeven van de stand van zaken binnen het onderzoeksdomein, op basis van een literatuurstudie.

In Hoofdstuk~\ref{ch:methodologie} wordt de methodologie toegelicht en worden de gebruikte onderzoekstechnieken besproken om een antwoord te kunnen formuleren op de onderzoeksvragen.

% TODO: Vul hier aan voor je eigen hoofstukken, één of twee zinnen per hoofdstuk

In Hoofdstuk~\ref{ch:conclusie}, tenslotte, wordt de conclusie gegeven en een antwoord geformuleerd op de onderzoeksvragen. Daarbij wordt ook een aanzet gegeven voor toekomstig onderzoek binnen dit domein.
\chapter{\IfLanguageName{dutch}{Stand van zaken}{State of the art}}%
\label{ch:stand-van-zaken}

% Tip: Begin elk hoofdstuk met een paragraaf inleiding die beschrijft hoe
% dit hoofdstuk past binnen het geheel van de bachelorproef. Geef in het
% bijzonder aan wat de link is met het vorige en volgende hoofdstuk.

% Pas na deze inleidende paragraaf komt de eerste sectiehoofding.

De beveiligingsrisico's van MCP-integraties in de financiële sector zijn pas goed te begrijpen tegen de achtergrond van hoe het protocol werkt en welke aanvalsvectoren het introduceert. Dit hoofdstuk brengt die context in kaart op basis van recente literatuur, en werkt toe naar de onderzoeksnood die de grondslag vormt voor de methodologie in Hoofdstuk~\ref{ch:methodologie}.

\section{Opkomst en risico's van MCP in Fintech}

De integratie van AI-agents in de financiële sector, voor toepassingen variërend van klantenservice tot transactieverwerking, leunt steeds meer op gestandaardiseerde protocollen. Het Model Context Protocol (MCP) verbindt AI-modellen met de bijbehorende bedrijfsapplicaties en databronnen. Hoewel deze standaardisatie integratie vergemakkelijkt, hebben \textcite{song2025protocolunveilingattackvectors} aangetoond dat het ecosysteem hierdoor nieuwe aanvalsvectoren introduceert die verder reiken dan het protocol zelf.

Die complexiteit brengt specifieke beveiligingsrisico's mee. Uit een systematische analyse door \textcite{guo2025systematicanalysismcpsecurity} blijkt dat MCP-implementaties vatbaar zijn voor diverse fundamentele kwetsbaarheden, waaronder \emph{prompt injection} via tool-beschrijvingen en \emph{privilege escalation}. Kwaadwillenden kunnen die zwakheden misbruiken. Zo beschrijven \textcite{Cato2025} in hun onderzoek naar "Living off AI"-aanvallen hoe MCP-servers gemanipuleerd kunnen worden omdat ze onvoldoende onderscheid maken tussen legitieme gebruikers en externe actoren, ook wel bekend als de \emph{confused deputy} aanval.

\section{Classificatie van kwetsbaarheden}

\textcite{OWASP2025} brachten de meest kritieke risico's in kaart in hun \emph{OWASP Top 10 for MCP}, met onder andere \emph{tool poisoning} en het gebrek aan degelijke audit logging als prominente bevindingen. Het is belangrijk om deze lijst te onderscheiden van de \emph{OWASP Top 10 for LLM Applications}~\autocite{OWASP2025}: de LLM-lijst richt zich op kwetsbaarheden in taalmodellen zelf (zoals onveilige output handling en training data poisoning), terwijl de MCP-lijst de beveiligingsrisico's van het protocol en de tool-interface beschrijft. Beide lijsten overlappen bij prompt injection en output sanitisatie, maar de MCP-specifieke risico's (zoals tool poisoning via kwaadaardige tool-beschrijvingen en het ontbreken van autorisatiechecks op tool-niveau) vallen buiten de scope van de LLM-lijst. In dit onderzoek hanteren we de \emph{OWASP Top 10 for MCP} als primaire taxonomie; waar relevant verwijzen we naar de LLM-lijst voor gedeelde risico's.

Voor fintech-bedrijven, die met gevoelige persoons- en transactiegegevens werken, vormen de MCP-specifieke kwetsbaarheden een direct risico op datalekken en compliance-schendingen. De OWASP MCP-taxonomie categoriseert aanvalsvectoren die algemene frameworks zoals de OWASP Top 10 voor webapplicaties niet dekken.

\textcite{hou2025mcplandscape} bieden een aanvullende taxonomie die MCP-dreigingen indeelt in vier aanvalstypen (kwaadwillende ontwikkelaars, externe aanvallers, misbruikende gebruikers en structurele implementatiefouten), verdeeld over 16 concrete bedreigingsscenario's. Het framework in dit onderzoek richt zich primair op de eerste en vierde categorie: kwetsbaarheden die een ontwikkelaar onbewust introduceert in tool-beschrijvingen en inputschema's. \textcite{huang2026mcpthreatmodeling} tonen via STRIDE-dreigings\-modellering aan dat de meeste MCP-clients geen statische validatie uitvoeren op ontvangen tool-metadata, wat precies de lacune is die de SAST-module in dit onderzoek opvult.

\section{Status Quo: Beperkingen van handmatige validatie}

Ondanks de ernst van de dreigingen, ontbreken er momenteel gestandaardiseerde methodologieën voor validatie. De huidige praktijk in het bedrijfsleven leunt sterk op handmatige en herhaalde code reviews en het gebruik van ad-hoc tools. \textcite{Uproot2025} bieden weliswaar handleidingen voor het uitvoeren van MCP-pentests, maar deze processen blijven arbeidsintensief en foutgevoelig. Een handmatige security review van een MCP-server implementatie kan gemiddeld 8 tot 16 uur in beslag nemen per iteratie, afhankelijk van de complexiteit van de server \autocite{Uproot2025}. In teams die frequent deployen, maakt dat grondige beveiligingscontroles moeilijk vol te houden.

Bestaande tools bieden slechts gedeeltelijke oplossingen. Hoewel er open-source initiatieven zijn, zoals de scanning tool van \textcite{Invariant2025}, ontbreekt vaak een integratie in een breder security framework dat geschikt is voor enterprise-omgevingen. Bedrijven vallen hierdoor terug op algemene tools zoals Burp Suite of Semgrep. Burp Suite kan JSON-RPC-verkeer onderscheppen maar mist de contextuele kennis om kwaadaardige tool-beschrijvingen te herkennen. Semgrep detecteert codepatronen maar kan de runtime-interactie tussen een agent en een MCP-server niet analyseren. Beide tools vullen daarmee elk een ander deel van het aanvalsvlak in, maar niet het MCP-specifieke deel.

Het ontbreken van zo'n framework is een onderbelicht gebied in zowel de literatuur als de praktijk, al evolueert het domein snel: de meeste relevante publicaties dateren uit 2025 en het MCP-ecosysteem is te jong voor definitieve generalisaties.

\section{Bestaande tools voor MCP-beveiligingstesting}

Er bestaan al enkele MCP-specifieke tools die gedeeltelijke dekking bieden. De meest directe is MCPSafetyScanner, ontwikkeld door \textcite{radosevich2025mcpsafetyaudit}. MCPSafetyScanner hanteert een agentische benadering: een LLM-agent genereert adversariale testscenario's en evalueert de respons van de doelserver. Dit maakt de tool contextbewust en flexibel in de keuze van testcases, maar ook niet-deterministisch: twee opeenvolgende runs op dezelfde server kunnen verschillende bevindingen opleveren, afhankelijk van het gedrag van het onderliggende model. Bovendien vereist de tool een actieve API-verbinding met een extern taalmodel, wat in gereguleerde fintech-omgevingen toestemmings- en dataprivacyvraagstukken meebrengt.

Het framework in dit onderzoek kiest bewust voor een deterministische aanpak. Patroonherkenning op basis van vaste detectieregels en vooraf gedefinieerde fuzzingpayloads vervangen de LLM-redenering van MCPSafetyScanner. Dit gaat ten koste van semantische flexibiliteit, maar garandeert reproduceerbaarheid (NF-2) en volledige onafhankelijkheid van externe diensten.

Een complementaire benadering biedt MCPTox~\autocite{wang2025mcptox}: een benchmark van 45 echte MCP-servers en 353 authentieke tools waarmee de robuustheid van \emph{LLM-agenten} tegen tool poisoning wordt gemeten. MCPTox rapporteert een aanvalssuccesratio van 72{,}8\% op o1-mini, wat de ernst van tool poisoning in productieomgevingen empirisch aantoont. Het onderscheid in evaluatieperspectief is hier van belang: MCPTox test of een model bestand is tegen een aanval, terwijl dit framework test of een server de aanvalsvector überhaupt \emph{bevat}. Beide invalshoeken vullen elkaar aan.

\section{Geautomatiseerde beveiligingstesting als oplossing}

De softwarebeveiligingsliteratuur beschrijft een reeks geautomatiseerde testmethoden die de nadelen van handmatige validatie kunnen opvangen. \emph{Static Application Security Testing} (SAST) analyseert broncode zonder uitvoering, wat het geschikt maakt voor vroege detectie van onveilige patronen in tool-beschrijvingen en configuratiebestanden \autocite{chess2004static}. \emph{Dynamic Application Security Testing} (DAST) test een draaiende applicatie door gemanipuleerde invoer te sturen en de respons te observeren. Deze aanpak is direct toepasbaar op de JSON-RPC-interface van MCP-servers. \emph{Fuzzing} ten slotte genereert automatisch grote volumes onverwachte of grensgevallen-invoer om onvoorziene kwetsbaarheden bloot te leggen \autocite{manes2019art}.

Toegepast op MCP vullen deze drie methoden elkaar aan: SAST detecteert statische patronen in tool-metadata en configuratiebestanden, terwijl DAST en fuzzing runtime-gedrag blootleggen dat niet uit de broncode alleen af te leiden is. Zo kunnen gemanipuleerde tool-beschrijvingen pas tijdens uitvoering hun effect tonen op het gedrag van een agent. Een framework dat deze drie benaderingen combineert, dekt zowel de statische als de dynamische aanvalsvlakken van MCP-servers.

\section{Kritische noot bij de literatuur}

De literatuur over MCP-beveiliging is jong: de meeste aangehaalde publicaties dateren uit 2025, en de OWASP MCP Top 10 werd pas gepubliceerd nadat MCP zelf breed werd geadopteerd. Dat heeft twee gevolgen voor de interpretatie. Sommige claims in de literatuur (waaronder de classificaties van \textcite{hou2025mcplandscape} en de dreigingsmodellen van \textcite{huang2026mcpthreatmodeling}) zijn nog niet empirisch gevalideerd op grote schaal. Daarnaast evolueert het MCP-ecosysteem snel: nieuwe protocol-versies, uitbreidingen en adoptie door grote platformen kunnen de relevantie van specifieke aanvalsvectoren snel doen verschuiven. Uitspraken over "het risico van MCP" moeten dan ook gelezen worden als een momentopname van de situatie in 2025, niet als definitieve conclusies.



%%=============================================================================
%% Methodologie
%%=============================================================================

\chapter{\IfLanguageName{dutch}{Methodologie}{Methodology}}%
\label{ch:methodologie}

%% TODO: In dit hoofstuk geef je een korte toelichting over hoe je te werk bent
%% gegaan. Verdeel je onderzoek in grote fasen, en licht in elke fase toe wat
%% de doelstelling was, welke deliverables daar uit gekomen zijn, en welke
%% onderzoeksmethoden je daarbij toegepast hebt. Verantwoord waarom je
%% op deze manier te werk gegaan bent.
%% 
%% Voorbeelden van zulke fasen zijn: literatuurstudie, opstellen van een
%% requirements-analyse, opstellen long-list (bij vergelijkende studie),
%% selectie van geschikte tools (bij vergelijkende studie, "short-list"),
%% opzetten testopstelling/PoC, uitvoeren testen en verzamelen
%% van resultaten, analyse van resultaten, ...
%%
%% !!!!! LET OP !!!!!
%%
%% Het is uitdrukkelijk NIET de bedoeling dat je het grootste deel van de corpus
%% van je bachelorproef in dit hoofstuk verwerkt! Dit hoofdstuk is eerder een
%% kort overzicht van je plan van aanpak.
%%
%% Maak voor elke fase (behalve het literatuuronderzoek) een NIEUW HOOFDSTUK aan
%% en geef het een gepaste titel.

Dit onderzoek hanteert een empirisch-technische aanpak op basis van experimentele validatie met kwantitatieve metrieken. Deze keuze past bij de onderzoeksvraag: om uitspraken te doen over detectienauwkeurigheid en tijdsbesparing zijn meetbare resultaten nodig die reproduceerbaar zijn. Het onderzoek verloopt in vier opeenvolgende fasen die op elkaar voortbouwen: een literatuurstudie legt het theoretische fundament, een requirementsanalyse vertaalt compliance-eisen naar concrete criteria, een prototypingfase realiseert het framework, en een validatiefase meet de effectiviteit ervan. Figuur~\ref{fig:methodologie_flow} geeft een schematisch overzicht van deze doorlopen fasen en de iteratieve feedbacklus tussen validatie en ontwerp.

% --- TIKZ FLOWCHART ---
\begin{figure*}[htbp]
    \centering
    \begin{tikzpicture}[node distance=1.5cm, scale=0.9, transform shape]
        % Stijlen definiëren
        \tikzstyle{startstop} = [rectangle, rounded corners, minimum width=3cm, minimum height=1cm, text centered, draw=black, fill=red!30]
        \tikzstyle{process} = [rectangle, minimum width=3cm, minimum height=1cm, text centered, draw=black, fill=orange!30]
        \tikzstyle{io} = [trapezium, trapezium left angle=70, trapezium right angle=110, minimum width=3cm, minimum height=1cm, text centered, draw=black, fill=blue!30]
        \tikzstyle{arrow} = [thick,->,>=stealth]

        % Knopen plaatsen
        \node (lit) [startstop] {1. Literatuurstudie \& Analyse};
        \node (req) [process, below of=lit] {2. Requirements Analyse (Fintech)};
        \node (design) [process, below of=req] {3. Ontwerp Framework \& Prototype};
        \node (poc) [process, below of=design] {4. Implementatie 'Goat' Server};
        \node (test) [process, below of=poc] {5. Validatie \& Benchmarking};
        \node (concl) [startstop, below of=test] {6. Conclusie \& Aanbevelingen};

        % Pijlen trekken
        \draw [arrow] (lit) -- (req);
        \draw [arrow] (req) -- (design);
        \draw [arrow] (design) -- (poc);
        \draw [arrow] (poc) -- (test);
        \draw [arrow] (test) -- (concl);

        % Feedback loop
        \draw [arrow] (test.east) -- ++(1.8,0) |- (design.east) node[pos=0.25, right, font=\small] {Iteratieve verbetering};

    \end{tikzpicture}
    \caption{Schematische weergave van de onderzoeksstappen.}
    \label{fig:methodologie_flow}
\end{figure*}

\section{Fase 1: Literatuurstudie en classificatie van dreigingen}

De eerste fase bestaat uit een literatuurstudie via academische databanken (Google Scholar, IEEE Xplore, arXiv) en officiële technische documentatie. De architectuur en datastromen van MCP worden doorgrond, en bestaande taxonomieën van AI-kwetsbaarheden --- zoals de OWASP Top 10 voor LLM's --- worden geanalyseerd op relevantie voor de MCP-context. Het concrete doel is een kwetsbaarheidstaxonomie opstellen voor prompt injection, tool poisoning en privilege escalation in financiële toepassingen. De resultaten vormen het hoofdstuk Stand van zaken. Geschatte doorlooptijd: vier weken.

\section{Fase 2: Requirementsanalyse voor Fintech-toepassingen}

De requirementsanalyse steunt op documentanalyse van compliance-kaders (PCI-DSS, GDPR) en API-security best practices, niet op interviews of enquêtes. Die keuze is bewust: harde compliance-eisen leveren objectievere criteria op dan subjectieve meningen. De functionele en niet-functionele vereisten --- waaronder logging, auditability en toegangscontrole --- worden ingedeeld via de MoSCoW-methode en gedocumenteerd in het hoofdstuk Requirementsanalyse. Geschatte doorlooptijd: twee weken.

\section{Fase 3: Ontwerp en implementatie van het assessment framework}

In de derde fase wordt het framework gebouwd als een modulaire Python-applicatie. Python is gekozen vanwege de native ondersteuning voor JSON-RPC en asynchrone communicatie, beide vereist voor MCP. Het framework bestaat uit drie kerncomponenten:

\begin{enumerate}
    \item \textbf{Discovery Module:} Analyseert de \texttt{list\_tools} en \texttt{list\_prompts} endpoints.
    \item \textbf{Static Analysis Engine:} Scant tool-beschrijvingen op onveilige patronen.
    \item \textbf{Dynamic Probing Engine:} Voert actieve fuzzing uit met gemanipuleerde payloads.
\end{enumerate}

De implementatie maakt gebruik van Python~3.12, Docker, Pydantic en de officiële \texttt{mcp}-library. Het resultaat is een werkende Proof of Concept gepubliceerd als open-source code op GitHub. Geschatte doorlooptijd: zes weken.

\section{Fase 4: Validatie via een 'Vulnerable MCP Server'}

Omdat er weinig publieke testdata beschikbaar is, wordt een eigen testomgeving ontwikkeld: de \emph{Fintech MCP Goat}, een opzettelijk kwetsbare MCP-server die financiële scenario's nabootst (saldo opvragen, overschrijving starten). Deze server dient als gecontroleerde \emph{ground truth}. Parallel wordt een gestructureerde handmatige beoordelingsprocedure vastgelegd als checklist, zodat de vergelijking reproduceerbaar is. Per testscenario worden vier grootheden gemeten:

\begin{itemize}
    \item \textbf{Detectieratio (Recall):} $\frac{\text{Aantal gevonden kwetsbaarheden}}{\text{Totaal aantal aanwezige kwetsbaarheden}}$, gemeten voor zowel de handmatige review als het framework
    \item \textbf{False Positive Rate:} Aantal onterechte alarmen voor beide methoden
    \item \textbf{Performance:} Benodigde tijd voor volledige assessment (handmatig in uren, geautomatiseerd in minuten)
    \item \textbf{Kwalitatieve analyse:} Identificatie waar menselijke expertise beter presteert (contextuele interpretatie) versus waar de tool superieur is (herhaalbaarheid, volledigheid)
\end{itemize}

De testomgeving draait via Docker~Compose; Burp~Suite wordt ingezet voor manuele verificatie. Het resultaat is een validatierapport met kwantitatieve meetresultaten. Geschatte doorlooptijd: drie weken.

\section{Overzichtstabel Planning}

Tabel~\ref{tab:planning} geeft een samenvattend overzicht van de fasering en deliverables.

\begin{table*}[htbp]
    \centering
    \caption{Planning en deliverables per onderzoeksfase.}
    \label{tab:planning}
    \begin{tabular}{@{}llp{6cm}@{}}
        \toprule
        \textbf{Fase}       & \textbf{Duur} & \textbf{Deliverable}               \\ \midrule
        1. Literatuurstudie & 4 weken       & Dreigingsmodel \& State-of-the-art \\
        2. Requirements     & 2 weken       & Requirementslijst (MoSCoW)         \\
        3. Implementatie    & 6 weken       & Python Framework (GitHub repo)     \\
        4. Validatie        & 3 weken       & Validatierapport \& Testresultaten \\
        5. Conclusie        & 2 weken       & Finaliseren Bachelorproef tekst    \\ \bottomrule
    \end{tabular}
\end{table*}



% Voeg hier je eigen hoofdstukken toe die de ``corpus'' van je bachelorproef
% vormen. De structuur en titels hangen af van je eigen onderzoek. Je kan bv.
% elke fase in je onderzoek in een apart hoofdstuk bespreken.

%\input{...}
%\input{...}
%...

%%=============================================================================
%% Conclusie
%%=============================================================================

\chapter{Conclusie}%
\label{ch:conclusie}

De hoofdvraag van dit onderzoek luidde: \emph{Hoe kan een framework worden ontwikkeld dat MCP-server implementaties automatisch test op een vooraf gedefinieerde set kritieke kwetsbaarheden, en hoe verhoudt de dekkingsgraad en benodigde tijd per assessment zich tot een handmatige beoordeling op basis van een concrete evaluatiecase?}

Het onderzoek toont aan dat een geautomatiseerd assessment framework kritieke kwetsbaarheden in MCP-serverimplementaties effectief kan detecteren. Het framework werkt sneller, levert reproduceerbare resultaten en detecteert bekende kwetsbaarheidspatronen even goed als een handmatige beoordeling.\footnote{Dekkingsgraad (recall): het aandeel daadwerkelijk aanwezige kwetsbaarheden dat het framework correct detecteert.} Voor contextuele interpretatie en complexe logische beveiligingsproblemen blijft menselijke analyse onmisbaar.

\paragraph{Kernresultaten}
Het framework detecteerde alle vijf ingebedde kwetsbaarheden in de Goat-server (zie \ref{sec:opstelling}) in minder dan 30 seconden, bij een precisie van 92{,}3\% in het meest conservatieve geval (12 terechte bevindingen op 13 totale bevindingen). De handmatige procedure vereiste 47 minuten en identificeerde 4 van de vijf kwetsbaarheden.

\section{Probleemdomein}

Het probleemdomein omvat de aard, de spreiding en de meetbaarheid van beveiligingskwetsbaarheden in MCP-servers; het oplossingsdomein (zie volgende sectie) omvat de detectiemechanismen van het framework: SAST-001 detecteert bijvoorbeeld de prompt-injection-patronen die in het probleemdomein zijn geïdentificeerd.

De literatuurstudie identificeerde prompt injection, tool poisoning en ontbrekende autorisatiechecks als de meest kritieke kwetsbaarheden voor Fintech-toepassingen. Deze drie categorieën scoren hoog op zowel kans als impact binnen de OWASP MCP Top 10, omdat ze direct raken aan bedrijfskritische processen zoals transactieverwerking en klantidentificatie.

Generieke beveiligingstools zoals Burp Suite en Semgrep dekken deze kwetsbaarheden slechts gedeeltelijk. Burp Suite test HTTP-verkeer maar kent de semantiek van MCP-schema's niet: het herkent geen kwaadaardige tool-beschrijvingen of poisoned responses. Semgrep kan patroonregels uitvoeren maar vereist MCP-specifieke rulesets die bij aanvang van dit onderzoek niet beschikbaar waren.

De handmatige procedure kostte 47 minuten. De voornaamste bronnen van foutmarge zijn aandachtsverlies bij repetitieve checks, inconsistente toepassing van criteria en de afhankelijkheid van de ervaring van de reviewer.

\section{Oplossingsdomein}

Het framework detecteerde 5 van de 5 kwetsbaarheden uit de ground truth (recall 100\%). Van de 13 bevindingen is er één debateerbaar als false positive, wat resulteert in een precisie van minimaal 92{,}3\% (12 terechte bevindingen op 13 totaal). Een klassieke false positive \emph{rate} (FPR = FP / (FP + TN)) is hier niet van toepassing omdat true negatives in een vulnerability scan niet zinvol te tellen zijn; de precisie is de correcte maatstaf. De handmatige beoordeling vond 4 van de vijf kwetsbaarheden (80\%) en produceerde 0 false positives.

Het framework presteert betrouwbaarder voor gedragsmatige kwetsbaarheden die alleen zichtbaar zijn bij runtime. Tool poisoning in \texttt{fetch\_market\_news} is enkel detecteerbaar door de tool effectief aan te roepen en de respons te analyseren; statische analyse alleen is hiervoor onvoldoende. Menselijke expertise blijft noodzakelijk voor de beoordeling van bedrijfsimpact: of een gevonden kwetsbaarheid in een specifieke architectuur exploiteerbaar is, vereist contextuele kennis die regels niet kunnen overnemen.

\section{Kritische reflectie}

De resultaten steunen op één gecontroleerde evaluatiecase (de Goat-server). De externe validiteit is daarmee beperkt: uitspraken over MCP-servers in het algemeen zijn op basis van dit onderzoek niet gerechtvaardigd. De Goat-server is opzettelijk kwetsbaar en representeert niet de gemiddelde productieomgeving.

Het hoge recall-cijfer is deels te verklaren door de aard van de testomgeving: de kwetsbaarheden zijn geselecteerd en geïmplementeerd om detecteerbaar te zijn voor dit type scanner. In een omgeving met subtielere of domeinspecifieke kwetsbaarheden zal de recall waarschijnlijk lager liggen.

Het debateerbare false positive op \texttt{list\_recent\_transactions} illustreert een inherent probleem bij patroondetectie. De scanner kan niet bepalen of een observatie een beveiligingsprobleem vormt zonder domeincontext. Dit vereist altijd een menselijke triage-stap na de geautomatiseerde scan.

\section{Aanbevelingen}

\paragraph{Gebruik het framework als screeningslaag, niet als auditvervanging}
De scanner identificeert patronen snel en reproduceerbaar, maar kan de bedrijfsimpact van een bevinding niet beoordelen. Of een gevonden autorisatieprobleem exploiteerbaar is binnen een specifieke architectuur, hangt af van netwerktopologie, toegangsbeheer en systeemprompts die buiten het bereik van de scanner liggen. Het debateerbare false positive op \texttt{list\_recent\_transactions} illustreert dat elke geautomatiseerde run gevolgd moet worden door een menselijke triagefase. Het framework vervangt geen security audit; het verlaagt de drempel om die audit gericht te voeren.

\paragraph{Verplicht autorisatiechecks op serverniveau}
\begin{sloppypar}
De kwetsbaarheden V2 en V3, het ontbreken van autorisatiecontroles op \texttt{analyze\_financial\_document} en \texttt{execute\_transfer}, vormen het hoogste risico in een Fintech-context.
\end{sloppypar} Beide tools gaven toegang tot gevoelige functionaliteit zonder enige verificatie van de aanroeper. Serverimplementaties dienen authenticatie en autorisatie af te dwingen op elke tool die toegang biedt tot transacties, gebruikersgegevens of beheersfuncties, ongeacht of de aanroeper een menselijke gebruiker of een AI-agent is.

\paragraph{Beschouw tool outputs als onvertrouwde invoer}
De tool poisoning in \texttt{fetch\_market\_news} toont aan dat kwetsbare inhoud ook via de teruggestuurde data kan binnenkomen, niet enkel via parameters. De verborgen injection-payload in de \texttt{content}-string was alleen detecteerbaar door de tool effectief aan te roepen en de response te analyseren. Servers dienen hun outputs te sanitiseren op dezelfde manier als inputs, met name wanneer de respons wordt doorgegeven aan een taalmodel dat de inhoud als instructie kan interpreteren.

\paragraph{Integreer de scanner in het CI/CD-proces}
Een volledige assessment duurt minder dan 30 seconden. Dit maakt continue beveiligingsmonitoring haalbaar: de scanner kan als stap in een bouwpijplijn worden opgenomen zodat kwetsbaarheden worden gedetecteerd vóór deployment. De gecontaineriseerde opzet via Docker Compose verlaagt de integratiedrempel, omdat de testomgeving identiek is aan de ontwikkelomgeving.

\section{Toekomstig onderzoek}

Drie vragen liggen open voor vervolgonderzoek. De externe validiteit vraagt om een uitgebreidere evaluatie: een studie met meerdere diverse MCP-servers, inclusief productieomgevingen, zou robuustere uitspraken mogelijk maken over de generaliseerbaarheid van de detectieregels.

Een logische vervolgstap is integratie in een CI/CD-pipeline. Het framework draait nu als zelfstandige tool; continuous security monitoring vereist dat de scanner ingebouwd wordt in het bouwproces, zodat bevindingen automatisch worden gerapporteerd aan het development team.

Een derde open vraag is hoe het framework zich gedraagt bij semantische kwetsbaarheden die afhangen van de systeemprompt van het onderliggende taalmodel. Of een prompt injection-patroon een specifiek model in de praktijk manipuleert, is niet deterministisch testbaar met de huidige aanpak. Dit vereist aanvullend onderzoek naar adversarial testing van LLM-agentsystemen.


%---------- Bijlagen -----------------------------------------------------------

\appendix

\chapter{Onderzoeksvoorstel}

Het onderwerp van deze bachelorproef is gebaseerd op een onderzoeksvoorstel dat vooraf werd beoordeeld door de promotor. Dat voorstel is opgenomen in deze bijlage.

%% TODO: 
\section*{Samenvatting}

% Kopieer en plak hier de samenvatting (abstract) van je onderzoeksvoorstel.
Binnen de Fintech-sector neemt de integratie van autonome AI-agents in bedrijfskritische processen structureel toe. Het Model Context Protocol (MCP) vormt hierbij de technische interface die AI-modellen verbindt met interne databronnen en API's. Hoewel MCP operationele voordelen biedt door interoperabiliteit te garanderen, is er momenteel een kennislacune over de effectiviteit van bestaande beveiligingsmethoden voor dit specifieke protocol. De centrale onderzoeksvraag luidt: hoe kan een framework worden ontwikkeld dat MCP-server implementaties automatisch test op een vooraf gedefinieerde set kritieke kwetsbaarheden, en hoe verhoudt de dekkingsgraad en benodigde tijd per assessment zich tot de huidige handmatige beveiligingsprocedures?
    
De doelstelling van dit onderzoek is drieledig: (1) het ontwikkelen van een werkend prototype dat minimaal de OWASP MCP Top 10 kwetsbaarheden kan detecteren, (2) een kwantitatieve validatie waarbij detectiegraad, tijdsbesparing en false positive/negative ratio's worden vergeleken met handmatige reviews, en (3) het formuleren van evidence-based aanbevelingen voor veilige MCP-implementaties. De methodologie combineert literatuurstudie, requirements-analyse op basis van compliance-eisen (PCI-DSS, GDPR), prototyping in Python, en experimentele validatie via een opzettelijk kwetsbare 'Fintech MCP Goat' testomgeving.
    
De verwachte resultaten omvatten een significant hogere en consistentere detectieratio voor technische kwetsbaarheden zoals prompt injection en tool poisoning, alsmede een substantiële tijdsbesparing bij het scannen van complexe MCP-servers. Voor de doelgroep — security teams, DevSecOps-engineers en enterprise architects — biedt dit onderzoek concrete meerwaarde door continuous compliance mogelijk te maken via CI/CD-integratie, de time-to-market te verbeteren door directe feedback aan ontwikkelaars, en het risico op fraude en datalekken aanzienlijk te reduceren.



% Verwijzing naar het bestand met de inhoud van het onderzoeksvoorstel
song2025protocolunveilingattackvectors%---------- Inleiding ---------------------------------------------------------

% TODO: Is dit voorstel gebaseerd op een paper van Research Methods die je
% vorig jaar hebt ingediend? Heb je daarbij eventueel samengewerkt met een
% andere student?
% Zo ja, haal dan de tekst hieronder uit commentaar en pas aan.

%\paragraph{Opmerking}

% Dit voorstel is gebaseerd op het onderzoeksvoorstel dat werd geschreven in het
% kader van het vak Research Methods dat ik (vorig/dit) academiejaar heb
% uitgewerkt (met medesturent VOORNAAM NAAM als mede-auteur).
% 

\section{Inleiding}%
\label{sec:inleiding}

De sector van financiële technologie, of Fintech, ontwikkelt zich razendsnel. Dat komt door de adoptie van autonome AI-agents. Deze slimme systemen doen niet alleen simpele klussen. Ze nemen complexere processen over. Van het afhandelen van klantvragen tot het verwerken van terugbetalingen en transacties. Voor effectieve communicatie met interne databases en API's is het Model Context Protocol essentieel. MCP zorgt voor interoperabiliteit als standaard.

Deze vooruitgang in technologie heeft een nadeel. Het verhoogt de veiligheidsrisico's aanzienlijk en die worden groter. Neem een Fintech-startup die een AI-agent via MCP toegang geeft tot betaalsystemen. Dat opent nieuwe wegen voor aanvallen. Kwaadwillenden kunnen prompt injection gebruiken om instructies van de agent te veranderen. Of tool poisoning om slechte commando's in data te verbergen. Gevolgen variëren van datalekken tot privilege escalation. Daarmee krijgt een agent bijvoorbeeld per ongeluk admin-rechten.\autocite{song2025protocolunveilingattackvectors}

Security-teams hebben vaak niet genoeg middelen om MCP-implementaties goed te beschermen. De praktijk hangt af van handmatige code-reviews en tijdelijke tools. Die tools passen niet bij de details van MCP. Dit maakt een knelpunt. Innovatie vertraagt door lange checks. Of servers gaan online zonder echte controle, wat vaker gebeurt. Onderzoek toont de urgentie. Zesenvijftig procent van de bedrijven met AI-agents meldt lekken in MCP-omgevingen, \autocite{guo2025systematicanalysismcpsecurity}. Dit voorstel pakt die kloof aan met automatisering en standaardisatie.




%\begin{itemize}
%  \item kaderen thema
%  \item de doelgroep
%  \item de probleemstelling en (centrale) onderzoeksvraag
%  \item de onderzoeksdoelstelling
%\end{itemize}
%
%Denk er aan: een typische bachelorproef is \textit{toegepast onderzoek}, wat betekent dat je start vanuit een concrete probleemsituatie in bedrijfscontext, een \textbf{casus}. Het is belangrijk om je onderwerp goed af te bakenen: je gaat voor die \textit{ene specifieke probleemsituatie} op zoek naar een goede oplossing, op basis van de huidige kennis in het vakgebied.
%
%De doelgroep moet ook concreet en duidelijk zijn, dus geen algemene of vaag gedefinieerde groepen zoals \emph{bedrijven}, \emph{developers}, \emph{Vlamingen}, enz. Je richt je in elk geval op it-professionals, een bachelorproef is geen populariserende tekst. Eén specifiek bedrijf (die te maken hebben met een concrete probleemsituatie) is dus beter dan \emph{bedrijven} in het algemeen.
%
%Formuleer duidelijk de onderzoeksvraag! De begeleiders lezen nog steeds te veel voorstellen waarin we geen onderzoeksvraag terugvinden.
%
%Schrijf ook iets over de doelstelling. Wat zie je als het concrete eindresultaat van je onderzoek, naast de uitgeschreven scriptie? Is het een proof-of-concept, een rapport met aanbevelingen, \ldots Met welk eindresultaat kan je je bachelorproef als een succes beschouwen?
%


\section{Probleemstelling}
\label{sec:probleemstelling}

Het kernprobleem is dat Fintech-bedrijven momenteel geen gestructureerde en geautomatiseerde manier hebben om MCP-servers op kritieke veiligheidskwetsbaarheden te controleren vóór live-gang.

\begin{itemize}
    \item \textbf{Huidige situatie (Status Quo):} Security teams vertrouwen op handmatige code reviews, wat een bottleneck vormt en vaak onvolledig is. Er wordt gebruikgemaakt van ad-hoc tools zoals Burp Suite of Semgrep, maar deze missen specifieke context voor MCP-gerelateerde aanvallen.
    \item \textbf{Specifieke risico's:} Zonder degelijke screening zijn servers kwetsbaar voor:
    \begin{itemize}
        \item \emph{Prompt injection attacks} via manipulatie van tool-beschrijvingen.
        \item \emph{Token theft \& privilege escalation}, waarbij agents te brede rechten krijgen.
        \item \emph{Tool poisoning}, waarbij kwaadaardige instructies in data worden verstopt.
        \item \emph{Confused deputy attacks}, omdat servers vaak geen onderscheid maken tussen gebruikers.
        \item Gebrek aan \emph{audit logging}, waardoor acties niet traceerbaar zijn.
    \end{itemize}
    \item \textbf{Doelgroep:} Het onderzoek richt zich op security teams, DevSecOps-engineers en enterprise architects binnen Fintech, e-commerce en andere gereguleerde industrieën.
\end{itemize}

\section{Doelstelling}
\label{sec:doelstelling}

Deze bachelorproef wil een geautomatiseerd security assessment framework ontwerpen, ontwikkelen en valideren. Het richt zich op MCP-servers. De resultaten zijn concreet voor MCP-servers.
\begin{enumerate}
    \item Een werkend prototype (Proof-of-Concept) van de security scanner.
    \item Een validatie van de effectiviteit van dit framework ten opzichte van handmatige reviews (detectiegraad en tijdsbesparing).
    \item Een set aanbevelingen voor de veilige implementatie van MCP in bedrijfsomgevingen.
\end{enumerate}

\section{Onderzoeksvragen}
\label{sec:onderzoeksvragen}

\paragraph{Hoofdvraag}
Hoe kan een geautomatiseerd security assessment framework ontworpen en geïmplementeerd worden om MCP-server implementaties in bedrijfsomgevingen systematischer en efficiënter te screenen op kritieke kwetsbaarheden dan de huidige handmatige processen?

\paragraph{Deelvragen Probleemdomein}
\begin{enumerate}
    \item Hoeveel van de top 10 MCP-kwetsbaarheden worden momenteel gedetecteerd door handmatige security reviews binnen fintech-bedrijven?
    \item Wat is het percentage MCP-implementaties in enterprises die kwetsbaar zijn voor prompt injection attacks, zoals vastgesteld door recent gepubliceerde onderzoeken?
    \item Hoe vaak leiden MCP-gerelateerde kwetsbaarheden in de praktijk tot datalekken of frauduleuze transacties (incidentratio in de laatste 24 maanden)?
    \item Hoeveel tijd besteden security teams gemiddeld aan het handmatig screenen van MCP-server code?
    \item Wat is het percentage false negatives bij bestaande handmatige controles?
\end{enumerate}

\paragraph{Deelvragen Oplossingsdomein}
\begin{enumerate}
    \item Hoeveel procent van de MCP-kwetsbaarheden kan een geautomatiseerde tool detecteren in vergelijking met handmatige reviews?
    \item Wat is de gemiddelde tijdsbesparing bij het scannen van MCP-servers door een geautomatiseerde pipeline?
    \item Wat is het percentage false positives/negatives bij gebruik van SAST, DAST en fuzzing op MCP-implementaties?
    \item Hoeveel MCP-servers kunnen effectief per dag worden gescand met het voorgestelde framework?
    \item Hoe verbeteren geprioriteerde rapporten de reactietijd van security teams?
\end{enumerate}


%---------- Stand van zaken ---------------------------------------------------

\section{Literatuurstudie}%
\label{sec:literatuurstudie}

\subsection{Opkomst en risico's van MCP in Fintech}

De integratie van AI-agents in de financiële sector, voor toepassingen variërend van klantenservice tot transactieverwerking, steunt in toenemende mate op gestandaardiseerde protocollen. Het Model Context Protocol (MCP) fungeert hierbij als een kritieke brug tussen AI-modellen en bedrijfsapplicaties met toebehorende data. Hoewel deze standaardisatie integratie vergemakkelijkt, hebben \textcite{song2025protocolunveilingattackvectors} aangetoond dat het ecosysteem hierdoor nieuwe aanvalsvectoren introduceert die verder reiken dan het protocol zelf.

De complexiteit van deze implementaties brengt specifieke beveiligingsrisico's met zich mee. Uit een systematische analyse door \textcite{guo2025systematicanalysismcpsecurity} blijkt dat MCP-implementaties vatbaar zijn voor diverse fundamentele kwetsbaarheden, waaronder \emph{prompt injection} via tool-beschrijvingen en \emph{privilege escalation}. In de praktijk kunnen kwaadwillenden deze zwakheden exploiteren. Zo beschrijven \textcite{Cato2025} in hun onderzoek naar "Living off AI"-aanvallen hoe MCP-servers gemanipuleerd kunnen worden omdat ze onvoldoende onderscheid maken tussen legitieme gebruikers en externe actoren, ook wel bekend als de \emph{confused deputy} aanval.

\subsection{Classificatie van kwetsbaarheden}

Om de grote hoeveelheid aan risico's te categoriseren, zijn er initiatieven gestart om de meest kritieke lekken in kaart te brengen. \textcite{OWASP2025} hebben een lijst van de "Top 10 MCP vulnerabilities" opgesteld, waarin onder andere \emph{tool poisoning} en het gebrek aan degelijke audit logging als prominente risico's worden bestempeld. Voor fintech-bedrijven, die met gevoelige persoons- en transactiegegevens werken, vormen deze specifieke kwetsbaarheden een direct risico op datalekken en compliance-schendingen.

\subsection{Status Quo: Beperkingen van handmatige validatie}

Ondanks de ernst van de dreigingen, ontbreken er momenteel gestandaardiseerde methodologieën voor validatie. De huidige praktijk in het bedrijfsleven leunt sterk op handmatige en herhaalde code reviews en het gebruik van ad-hoc tools. \textcite{Uproot2025} bieden weliswaar handleidingen voor het uitvoeren van MCP-pentests, maar deze processen blijven arbeidsintensief en foutgevoelig, wat het tijdsaspect dus niet positief beïnvloedt.

Bestaande tools bieden slechts gedeeltelijke oplossingen. Hoewel er open-source initiatieven zijn, zoals de scanning tool van \textcite{Invariant2025}, ontbreekt vaak een integratie in een breder security framework dat geschikt is voor enterprise-omgevingen. Bedrijven vallen hierdoor vaak terug op algemene tools zoals Burp Suite of Semgrep die de context-specifieke logica van MCP niet volledig begrijpen. Dit leidt tot een situatie waarin, volgens marktobservaties, een aanzienlijk deel van de deployments live gaat met niet-gedetecteerde veiligheidslekken.

Het ontbreken van een geautomatiseerd, gestructureerd framework dat specifiek is ontworpen voor de nuances van MCP-servers, vormt het centrale onderwerp in de huidige literatuur en praktijk. Dit onderzoek beoogt dit gat te dichten door een dergelijk framework te ontwerpen en de effectiviteit hiervan te evalueren.


%Hier beschrijf je de \emph{state-of-the-art} rondom je gekozen onderzoeksdomein, d.w.z.\ een inleidende, doorlopende tekst over het onderzoeksdomein van je bachelorproef. Je steunt daarbij heel sterk op de professionele \emph{vakliteratuur}, en niet zozeer op populariserende teksten voor een breed publiek. Wat is de huidige stand van zaken in dit domein, en wat zijn nog eventuele open vragen (die misschien de aanleiding waren tot je onderzoeksvraag!)?
%
%Je mag de titel van deze sectie ook aanpassen (literatuurstudie, stand van zaken, enz.). Zijn er al gelijkaardige onderzoeken gevoerd? Wat concluderen ze? Wat is het verschil met jouw onderzoek?
%
%Verwijs bij elke introductie van een term of bewering over het domein naar de vakliteratuur, bijvoorbeeld~\autocite{Hykes2013}! Denk zeker goed na welke werken je refereert en waarom.
%
%Draag zorg voor correcte literatuurverwijzingen! Een bronvermelding hoort thuis \emph{binnen} de zin waar je je op die bron baseert, dus niet er buiten! Maak meteen een verwijzing als je gebruik maakt van een bron. Doe dit dus \emph{niet} aan het einde van een lange paragraaf. Baseer nooit teveel aansluitende tekst op eenzelfde bron.
%
%Als je informatie over bronnen verzamelt in JabRef, zorg er dan voor dat alle nodige info aanwezig is om de bron terug te vinden (zoals uitvoerig besproken in de lessen Research Methods).
%
% Voor literatuurverwijzingen zijn er twee belangrijke commando's:
% \autocite{KEY} => (Auteur, jaartal) Gebruik dit als de naam van de auteur
%   geen onderdeel is van de zin.
% \textcite{KEY} => Auteur (jaartal)  Gebruik dit als de auteursnaam wel een
%   functie heeft in de zin (bv. ``Uit onderzoek door Doll & Hill (1954) bleek
%   ...'')

Je mag deze sectie nog verder onderverdelen in subsecties als dit de structuur van de tekst kan verduidelijken.

%---------- Methodologie ------------------------------------------------------
\section{Methodologie}%
\label{sec:methodologie}

Om objectieve en meetbare resultaten te garanderen, worden subjectieve methoden zoals enquêtes vermeden. In plaats daarvan wordt gekozen voor experimentele validatie aan de hand van kwantitatieve metrieken. Figuur~\ref{fig:methodologie_flow} geeft een schematisch overzicht van de doorlopen fasen.

% --- TIKZ FLOWCHART ---
\begin{figure*}[htbp]
\centering
\begin{tikzpicture}[node distance=1.5cm]
    % Stijlen definiëren
    \tikzstyle{startstop} = [rectangle, rounded corners, minimum width=3cm, minimum height=1cm, text centered, draw=black, fill=red!30]
    \tikzstyle{process} = [rectangle, minimum width=3cm, minimum height=1cm, text centered, draw=black, fill=orange!30]
    \tikzstyle{io} = [trapezium, trapezium left angle=70, trapezium right angle=110, minimum width=3cm, minimum height=1cm, text centered, draw=black, fill=blue!30]
    \tikzstyle{arrow} = [thick,->,>=stealth]

    % Knopen plaatsen
    \node (lit) [startstop] {1. Literatuurstudie \& Analyse};
    \node (req) [process, below of=lit] {2. Requirements Analyse (Fintech)};
    \node (design) [process, below of=req] {3. Ontwerp Framework \& Prototype};
    \node (poc) [process, below of=design] {4. Implementatie 'Goat' Server};
    \node (test) [process, below of=poc] {5. Validatie \& Benchmarking};
    \node (concl) [startstop, below of=test] {6. Conclusie \& Aanbevelingen};

    % Pijlen trekken
    \draw [arrow] (lit) -- (req);
    \draw [arrow] (req) -- (design);
    \draw [arrow] (design) -- (poc);
    \draw [arrow] (poc) -- (test);
    \draw [arrow] (test) -- (concl);
    
    % Feedback loop
    \draw [arrow] (test.east) -- ++(1,0) |- (design.east) node[pos=0.25, right] {Iteratieve verbetering};

\end{tikzpicture}
\caption{Schematische weergave van de onderzoeksstappen.}
\label{fig:methodologie_flow}
\end{figure*}

\subsection{Fase 1: Literatuurstudie en classificatie van dreigingen}

\textbf{Onderzoekstechniek:} Literatuurstudie. \\
Het onderzoek start met een diepgaande analyse van de technische specificaties van het Model Context Protocol (MCP). Er wordt gebruikgemaakt van academische databanken (zoals Google Scholar, IEEE Xplore en arXiv) en officiële technische documentatie om de architectuur en datastromen te doorgronden. Specifiek wordt gezocht naar bestaande taxonomieën van AI-kwetsbaarheden (zoals de OWASP Top 10 voor LLM's) en hoe deze zich vertalen naar de MCP-context.

\textbf{Doel:} Het opstellen van een specifieke taxonomie van kwetsbaarheden (zoals prompt injection, tool poisoning en privilege escalation) die relevant zijn voor de financiële sector. 

\textbf{Tools:} BiBTex voor referentiebeheer, LaTeX voor documentatie. \\
\textbf{Tijdschatting:} 4 weken. \\
\textbf{Deliverable:} Hoofdstuk 'Literatuurstudie' met een gedefinieerd dreigingsmodel.

\subsection{Fase 2: Requirementsanalyse voor Fintech-toepassingen}

\textbf{Onderzoekstechniek:} Documentanalyse en functionele analyse. \\
In plaats van interviews of enquêtes, die subjectieve meningen opleveren, wordt er gekeken naar harde compliance-eisen uit de financiële sector (zoals PCI-DSS en GDPR) en best practices voor API-security. Op basis hiervan worden de functionele en niet-functionele vereisten (zoals logging, auditability en toegangscontrole) voor de scanner vastgelegd volgens de MoSCoW-methode.

\textbf{Doel:} Een objectieve lijst van criteria opstellen waaraan het te ontwikkelen framework moet voldoen om inzetbaar te zijn in een enterprise-omgeving.

\textbf{Tools:} Requirements management templates. \\
\textbf{Tijdschatting:} 2 weken. \\
\textbf{Deliverable:} Een geprioriteerde lijst van requirements.

\subsection{Fase 3: Ontwerp en implementatie van het assessment framework}

\textbf{Onderzoekstechniek:} Prototyping en iteratieve softwareontwikkeling (PoC). \\
Dit is de technisch meest complexe fase waarin het artefact daadwerkelijk wordt gebouwd. De scanner wordt modulair opgezet om flexibiliteit te garanderen. De technische realisatie gebeurt in Python vanwege de uitgebreide ondersteuning voor JSON-RPC en asynchrone communicatie, essentieel voor MCP.

Het framework bestaat uit drie kerncomponenten:
\begin{enumerate}
    \item \textbf{Discovery Module:} Analyseert de \texttt{list\_tools} en \texttt{list\_prompts} endpoints.
    \item \textbf{Static Analysis Engine:} Scant tool-beschrijvingen op onveilige patronen.
    \item \textbf{Dynamic Probing Engine:} Voert actieve fuzzing uit met gemanipuleerde payloads.
\end{enumerate}

\textbf{Tools:} Python 3.12+, Docker (voor containerisatie), Git (versiebeheer), Pydantic (data validatie), `mcp`-library. \\
\textbf{Tijdschatting:} 6 weken. \\
\textbf{Deliverable:} Een werkende Proof of Concept (PoC) van de scanner, gepubliceerd als open-source code op GitHub.

\subsection{Fase 4: Validatie via een 'Vulnerable MCP Server'}

\textbf{Onderzoekstechniek:} Experimentele validatie en vergelijkende studie. \\
Omdat er weinig publieke testdata is, wordt een eigen testomgeving ontwikkeld: de 'Fintech MCP Goat'. Dit is een opzettelijk onveilige MCP-server die financiële use-cases nabootst (bv. saldo opvragen, overschrijving starten). 

Deze omgeving dient als 'ground truth'. De effectiviteit van de scanner wordt gemeten door deze los te laten op de 'Goat' server en de resultaten te vergelijken met de bekende kwetsbaarheden. De validatie is puur kwantitatief:
\begin{itemize}
    \item \textbf{Detectieratio (Recall):} $\frac{\text{Aantal gevonden kwetsbaarheden}}{\text{Totaal aantal aanwezige kwetsbaarheden}}$
    \item \textbf{False Positive Rate:} Aantal onterechte alarmen.
    \item \textbf{Performance:} Scan-tijd in milliseconden.
\end{itemize}

\textbf{Tools:} Docker Compose (voor het opzetten van de testlab), Pytest (voor geautomatiseerde testsuites), Burp Suite (voor manuele verificatie van bevindingen). \\
\textbf{Tijdschatting:} 3 weken. \\
\textbf{Deliverable:} Validatierapport met meetresultaten en een analyse van de nauwkeurigheid.

\subsection{Overzichtstabel Planning}

Tabel~\ref{tab:planning} geeft een samenvattend overzicht van de fasering en deliverables.

\begin{table*}[htbp]
\centering
\caption{Planning en deliverables per onderzoeksfase.}
\label{tab:planning}
\begin{tabular}{@{}llp{6cm}@{}}
\toprule
\textbf{Fase} & \textbf{Duur} & \textbf{Deliverable} \\ \midrule
1. Literatuurstudie & 4 weken & Dreigingsmodel \& State-of-the-art \\
2. Requirements & 2 weken & Requirementslijst (MoSCoW) \\
3. Implementatie & 6 weken & Python Framework (GitHub repo) \\
4. Validatie & 3 weken & Validatierapport \& Testresultaten \\
5. Conclusie & 2 weken & Finaliseren Bachelorproef tekst \\ \bottomrule
\end{tabular}
\end{table*}

%Hier beschrijf je hoe je van plan bent het onderzoek te voeren. Welke onderzoekstechniek ga je toepassen om elk van je onderzoeksvragen te beantwoorden? Gebruik je hiervoor literatuurstudie, interviews met belanghebbenden (bv.~voor requirements-analyse), experimenten, simulaties, vergelijkende studie, risico-analyse, PoC, \ldots?
%
%Valt je onderwerp onder één van de typische soorten bachelorproeven die besproken zijn in de lessen Research Methods (bv.\ vergelijkende studie of risico-analyse)? Zorg er dan ook voor dat we duidelijk de verschillende stappen terug vinden die we verwachten in dit soort onderzoek!
%
%Vermijd onderzoekstechnieken die geen objectieve, meetbare resultaten kunnen opleveren. Enquêtes, bijvoorbeeld, zijn voor een bachelorproef informatica meestal \textbf{niet geschikt}. De antwoorden zijn eerder meningen dan feiten en in de praktijk blijkt het ook bijzonder moeilijk om voldoende respondenten te vinden. Studenten die een enquête willen voeren, hebben meestal ook geen goede definitie van de populatie, waardoor ook niet kan aangetoond worden dat eventuele resultaten representatief zijn.
%
%Uit dit onderdeel moet duidelijk naar voor komen dat je bachelorproef ook technisch voldoen\-de diepgang zal bevatten. Het zou niet kloppen als een bachelorproef informatica ook door bv.\ een student marketing zou kunnen uitgevoerd worden.
%
%Je beschrijft ook al welke tools (hardware, software, diensten, \ldots) je denkt hiervoor te gebruiken of te ontwikkelen.
%
%Probeer ook een tijdschatting te maken. Hoe lang zal je met elke fase van je onderzoek bezig zijn en wat zijn de concrete \emph{deliverables} in elke fase?
%


%---------- Verwachte resultaten ----------------------------------------------
\section{Verwacht resultaat, conclusie}%
\label{sec:verwachte_resultaten}

\subsection{Verwachte prestaties van het framework}

Het belanrijkste verwachte resultaat is een functionerend prototype (de scanner) dat in staat is om de in de literatuurstudie geïdentificeerde top-kwetsbaarheden autonoom te detecteren. Er worden twee specifieke meetbare uitkomsten verwacht: een verhoogde detectieconsistentie en een significante tijdsbesparing.

\subsubsection{Hypothese 1: Detectieratio en Dekking}
De verwachting is dat het geautomatiseerde framework een hogere en consistentere detectieratio (Recall) behaalt voor technische kwetsbaarheden (zoals \emph{Prompt Injection} en onveilige \emph{Tool Definitions}) dan een gemiddelde handmatige review.

Figuur~\ref{fig:mockup_detectie} toont de verwachte verdeling van de detectieresultaten. Op de Y-as staat het percentage gedetecteerde kwetsbaarheden (Detectieratio), en op de X-as de verschillende categorieën van MCP-risico's.

\begin{figure*}[htbp]
    \centering
    \begin{tikzpicture}
        % Assen
        \draw[thick,->] (0,0) -- (8,0) node[anchor=north west] {Type Kwetsbaarheid};
        \draw[thick,->] (0,0) -- (0,5) node[anchor=south east] {Detectieratio (\%)};
        
        % Grid en labels Y-as
        \foreach \y/\label in {1/20, 2/40, 3/60, 4/80} {
            \draw[gray!30] (0,\y) -- (8,\y);
            \node[left] at (0,\y) {\label\%};
        }
        
        % Labels X-as
        \node[below, align=center] at (1.75,0) {\small Prompt\\ \small Injection};
        \node[below, align=center] at (4.25,0) {\small Tool\\ \small Poisoning};
        \node[below, align=center] at (6.75,0) {\small Privilege\\ \small Escalation};

        % Data (Mock-up) - Blauw = Framework, Grijs = Handmatig
        % Prompt Injection
        \fill[blue!60] (1.0,0) rectangle (1.5, 4.2); % Framework (hoog)
        \fill[gray!60] (2.0,0) rectangle (2.5, 2.5); % Handmatig (lager/wisselend)
        
        % Tool Poisoning
        \fill[blue!60] (3.5,0) rectangle (4.0, 3.8); 
        \fill[gray!60] (4.5,0) rectangle (5.0, 1.5); % Handmatig mist dit vaak
        
        % Priv Escalation
        \fill[blue!60] (6.0,0) rectangle (6.5, 3.5); 
        \fill[gray!60] (7.0,0) rectangle (7.5, 3.0); 

        % Legenda
        \fill[blue!60] (5,4.5) rectangle (5.3,4.8) node[right, black] {\small Framework};
        \fill[gray!60] (5,4.0) rectangle (5.3,4.3) node[right, black] {\small Handmatige Review};
    \end{tikzpicture}
    \caption{Mock-up: Verwachte detectieratio van het framework t.o.v. handmatige review. We verwachten dat het framework vooral beter scoort op complexe, technische patronen zoals Tool Poisoning.}
    \label{fig:mockup_detectie}
\end{figure*}

De motivatie achter deze verwachting is dat statische analyse-engines (SAST) in staat zijn om duizenden regels code en configuratiebestanden te scannen op patronen die voor het menselijk oog moeilijk te vinden zijn, zeker wanneer deze verspreid zijn over meerdere bestanden. Wel wordt verwacht dat het aantal \emph{False Positives} (onterechte alarmen) bij het framework initieel hoger zal liggen dan bij een menselijke expert, wat in de discussie genuanceerd zal moeten worden.

\subsubsection{Hypothese 2: Tijdsbesparing en Schaalbaarheid}
Een tweede belangrijke verwachting is de schaalbaarheid. Figuur~\ref{fig:mockup_tijd} visualiseert de verwachte tijdsbesteding. Op de X-as staat de complexiteit van de MCP-server (aantal tools/endpoints), en op de Y-as de benodigde tijd voor een security assessment.

\begin{figure*}[htbp]
    \centering
    \begin{tikzpicture}
        % Assen
        \draw[thick,->] (0,0) -- (7,0) node[anchor=north west] {Aantal MCP Tools/Endpoints};
        \draw[thick,->] (0,0) -- (0,5) node[anchor=south east] {Tijd (uren)};
        
        % Lijnen tekenen
        % Handmatig (Lineair steil)
        \draw[red, thick] (0,0.5) -- (6,4.5) node[right] {Handmatig};
        % Framework (Vlakker)
        \draw[blue, thick] (0,0.2) -- (6,1.2) node[right] {Framework};
        
        % Stippellijnen voor threshold
        \draw[dashed, gray] (3,0) -- (3,5);
        \node[above, gray] at (3,5) {Omslagpunt complexiteit};

    \end{tikzpicture}
    \caption{Mock-up: Verwachte tijdsbesteding. Naarmate de complexiteit van de MCP-server toeneemt, blijft de scantijd van het framework nagenoeg constant, terwijl handmatige review lineair toeneemt.}
    \label{fig:mockup_tijd}
\end{figure*}

De verwachting is dat het framework een volledige scan kan uitvoeren in enkele minuten, ongeacht of de server 5 of 50 tools definieert. Dit staat in contrast met handmatige penetratietesten, die dagen in beslag kunnen nemen.

\subsection{Meerwaarde voor de doelgroep}

Voor Fintech-bedrijven, de primaire doelgroep, biedt dit resultaat drie concrete voordelen:
\begin{enumerate}
    \item \textbf{Continuous Compliance:} Door de scanner te integreren in de CI/CD-pijplijn, kunnen bedrijven bij elke code-wijziging aantonen dat de AI-agents nog steeds voldoen aan veiligheidseisen. Dit is cruciaal voor audits (bv. SOC2 of DORA).
    \item \textbf{Time-to-Market:} Security wordt geen bottleneck meer. Ontwikkelaars krijgen directe feedback over onveilige tool-definities zonder te moeten wachten op een beschikbaar security-team.
    \item \textbf{Risicoreductie:} Door veelvoorkomende kwetsbaarheden zoals \emph{Confused Deputy} aanvallen automatisch af te vangen, verkleint de kans op fraude en datalekken via AI-interfaces aanzienlijk.
\end{enumerate}

\subsection{Conclusie}
Dit onderzoek toont aan dat een gespecialiseerd security framework voor MCP haalbaar en nodig is. Voor veilige AI-agents in financiën.

Als resultaten afwijken, zoals te weinig context voor complexe fouten, is dat ook waardevol. Het wijst op afhankelijkheid van menselijke expertise. En noodzaak voor hybrid approach in de nabije toekomst.



%Hier beschrijf je welke resultaten je verwacht. Als je metingen en simulaties uitvoert, kan je hier al mock-ups maken van de grafieken samen met de verwachte conclusies. Benoem zeker al je assen en de onderdelen van de grafiek die je gaat gebruiken. Dit zorgt ervoor dat je concreet weet welk soort data je moet verzamelen en hoe je die moet meten.
%
%Wat heeft de doelgroep van je onderzoek aan het resultaat? Op welke manier zorgt jouw bachelorproef voor een meerwaarde?
%
%Hier beschrijf je wat je verwacht uit je onderzoek, met de motivatie waarom. Het is \textbf{niet} erg indien uit je onderzoek andere resultaten en conclusies vloeien dan dat je hier beschrijft: het is dan juist interessant om te onderzoeken waarom jouw hypothesen niet overeenkomen met de resultaten.
%


%%---------- Andere bijlagen --------------------------------------------------
% TODO: Voeg hier eventuele andere bijlagen toe. Bv. als je deze BP voor de
% tweede keer indient, een overzicht van de verbeteringen t.o.v. het origineel.
%\input{...}

%%---------- Backmatter, referentielijst ---------------------------------------

\backmatter{}

\setlength\bibitemsep{2pt} %% Add Some space between the bibliograpy entries
\printbibliography[heading=bibintoc]

\end{document}
